\section{Spectral references and disconnected components}
\label{sec:spectral}

Phase is not the only datum that normalized operational records can forget.
The maximal mismatch rule of Section~\ref{sec:refinement} assigns a positive rate
$R=\lambda_{\max}(X)I-X$ to a Hermitian mismatch $X$.
Given a positive operator $H$, the choice $X=\lambda_{\max}(H)I-H$ yields
\begin{equation}
 R=C(H):=H-\lambda_{\min}(H)I.
 \label{eq:centering}
\end{equation}
This is invariant under every admissible shift $H\mapsto H+cI$.
Here $H$ has rate units, or the readout parameter has been calibrated to supply
the reciprocal units; the notation alone supplies no physical clock.
The no-jump amplitude $e^{-tR/2}$ and its survival probabilities therefore cannot
identify the absolute spectral floor.

\begin{proposition}[Disconnected-spectator obstruction]
\label{prop:centering}
For positive matrices $H_1,H_2$ of nonzero dimensions, let
$\lambda_i=\lambda_{\min}(H_i)$.
Then
\begin{equation}
 C(H_1\oplus H_2)|_i
 =C(H_i)+\bigl(\lambda_i-\min\{\lambda_1,\lambda_2\}\bigr)I_i.
 \label{eq:centering-defect}
\end{equation}
Consequently $C$ preserves this direct sum if and only if
$\lambda_1=\lambda_2$.
\end{proposition}
\begin{proof}
The spectrum of a direct sum is the union of its component spectra, so its
floor is the minimum of their floors.
Subtracting that common floor from each block gives
\eqref{eq:centering-defect}; both offsets vanish precisely in the stated case.
\end{proof}

For example, take
\begin{equation}
 H_1=\operatorname{diag}(1/5,11/10),\qquad
 H_2=\operatorname{diag}(7/10,3/2).
\end{equation}
The two results are
\begin{align}
 C(H_1)\oplus C(H_2)&=\operatorname{diag}(0,9/10,0,4/5),\nonumber\\
 C(H_1\oplus H_2)&=\operatorname{diag}(0,9/10,1/2,13/10).
\end{align}
The disconnected lower-floor component adds $I/2$ to the rate on the second
component.
At a fixed readout time, its survival probability is multiplied by $e^{-t/2}$.
Thus this defect is observable in a resolved no-jump record; it is not merely
an overall Hamiltonian phase.
Direct-sum preservation is the precise disconnected-component requirement here.
It should not be equated with every notion of relativistic locality or with
tensor-product composition.

\subsection{A general finite reference theorem}

Changing the scalar statistic in \eqref{eq:centering} cannot preserve all of the
same requirements while retaining spectral information.
The following statement makes this obstruction independent of the choice of
the minimum eigenvalue.

\begin{theorem}[Spectral-reference obstruction]
\label{thm:reference}
For every $n\geq1$, let
$F_n:\Herm(n)_+\to\Herm(n)$.
Assume unitary covariance, preservation of all finite direct sums, and scalar
shift invariance:
\begin{align}
 F_n(UHU^\dagger)&=UF_n(H)U^\dagger,
 \label{eq:reference-unitary}\\
 F_{n+m}(H\oplus K)&=F_n(H)\oplus F_m(K),
 \label{eq:reference-sum}\\
 F_n(H+cI)&=F_n(H)
 \label{eq:reference-shift}
\end{align}
whenever the arguments are positive semidefinite.
There is a real constant $a$, independent of $n,H$, such that
$F_n(H)=aI_n$.
If $F_1(0)=0$, every $F_n$ is zero.
No continuity, linearity, differentiability or spectral nondegeneracy
assumption is required.
\end{theorem}
\begin{proof}
Write $F_1([\lambda])=[f(\lambda)]$.
Equation~\eqref{eq:reference-shift} applied to $[0]$ and $[\lambda]$ gives
$f(\lambda)=f(0)=a$ for every $\lambda\geq0$.
By repeated use of \eqref{eq:reference-sum},
\[
 F_n(\operatorname{diag}(\lambda_1,\ldots,\lambda_n))
 =\operatorname{diag}(f(\lambda_1),\ldots,f(\lambda_n))=aI_n.
\]
The spectral theorem and \eqref{eq:reference-unitary} extend this equality to
every positive Hermitian $H$.
The final assertion follows from $a=F_1(0)$.
\end{proof}

The theorem is deliberately a statement about bare matrices in all finite
dimensions.
A model equipped with a preferred component algebra has additional data and
need not be covariant under every unitary that mixes those components.
Componentwise centering can then be defined relative to that algebra, but the
component choice has not been obtained from a shift-blind matrix record.
Replacing the minimum by a mean or quantile does not evade
Theorem~\ref{thm:reference} under its hypotheses.

\subsection{Calibrated repairs and a complete positive process}

There are two elementary repairs.
One may retain $H$ itself and its absolute spectral floor, thereby preserving
unitary covariance and direct sums while abandoning shift blindness.
Alternatively supply a reference operator $B$ with $R=H-B\geq0$.
The assignment
\begin{equation}
 R(H,B)=H-B,\qquad
 R(H+cI,B+cI)=R(H,B)
 \label{eq:calibrated-reference}
\end{equation}
preserves direct sums of pairs and simultaneous unitary transformations.
It contains an extra calibrated datum; it does not extract an absolute
generator from the original normalized record.

A complete process can be attached to any chosen $R\geq0$ without imposing a
two-Kraus chain law on a general spectrum.
Let $R=\sum_{j=1}^d r_j|j\rangle\langle j|$, introduce an absorbing flag
$|\varnothing\rangle$ orthogonal to the live space, and define
\begin{align}
 A_0(t)&=|\varnothing\rangle\langle\varnothing|
       +\sum_j e^{-tr_j/2}|j\rangle\langle j|,\nonumber\\
 A_j(t)&=\sqrt{1-e^{-tr_j}}\,
           |\varnothing\rangle\langle j|,\qquad j=1,\ldots,d .
 \label{eq:absorbing-completion}
\end{align}

\begin{proposition}[Calibrated absorbing completion]
\label{prop:absorbing}
The operators \eqref{eq:absorbing-completion} define a trace-preserving
completely positive semigroup on the live space plus the absorbing flag.
For a normalized live input $\rho$, its probability of remaining live is
$\Tr(e^{-tR}\rho)$.
For the maximally mixed live input this is $d^{-1}\Tr e^{-tR}$.
\end{proposition}
\begin{proof}
The completeness identity follows from
$e^{-tr_j}+(1-e^{-tr_j})=1$ for each live mode and the unit coefficient on the
absorbing flag.
Live matrix elements evolve by
$\rho_{jk}\mapsto e^{-t(r_j+r_k)/2}\rho_{jk}$, live--flag coherences by
$\rho_{j\varnothing}\mapsto e^{-tr_j/2}\rho_{j\varnothing}$, and the flag
population receives $\sum_j(1-e^{-tr_j})\rho_{jj}$.
These formulas compose at times $t_1,t_2$ by addition of the times.
Equivalently, the generator has jump operators
$L_j=\sqrt{r_j}|\varnothing\rangle\langle j|$.
Taking the trace of the live block gives the asserted probabilities.
\end{proof}

This construction is a calibrated measurement of a heat trace.
For a geometric choice $H=D^\dagger D$, convergence to a continuum spectrum
and the physical meaning of the readout time remain separate questions.
A protected zero mode anchors the floor only in sectors that contain it.
It is not universal: on the round $S^4_r$ the untwisted self-adjoint Dirac
operator has eigenvalues $\pm(k+2)/r$, so $D^2$ has floor $4/r^2$
\cite{Camporesi}.
The floor belongs to $D^2$, not to the two-sided Dirac spectrum itself.
Neither an absorbing completion nor a chosen reference $B$ supplies a
gravitational action or selects a microscopic clock.
